\section{Computational teleology and cybernetics: an open-ended universe near the undecidability boundary}
\label{sec:teleology}

\subsection{Teleology as a computable objective: sustaining open-endedness}
Once undecidable reachability predicates exist (Theorem~\ref{thm:undecidable_reachability}), the phrase ``the universe aims to compute an ultimate answer and halt'' is not the natural fixed point.
A more precise statement is available: for finite-resource observers, what matters is whether the scan--readout protocol continues to generate histories that are simultaneously (i) not trapped in short periodic loops (overly decidable/compressible) and (ii) not washed out into pure noise (unreadable/incompressible).

We package this as a definition at the interpretation/interface layer.

\begin{definition}[Computational teleology (interface definition)]
\label{def:computational_teleology}
A unitary scan--readout universe exhibits \emph{computational teleology} if its effective operation stabilizes a regime in which
\begin{itemize}
    \item \textbf{Expressivity:} local reachability remains rich enough to support undecidable predicates (open-endedness);
    \item \textbf{Legibility:} finite-resource observers can extract stable, compressible, and actionable readout histories.
\end{itemize}
Equivalently, teleology is an operational balance between undecidability (expressivity) and readout feasibility (resource-bounded distinguishability).
\end{definition}

\subsection{Golden anti-resonance as a control law}
The golden branch supplies a particularly rigid ``control'' mechanism: it avoids phase lock-in (anti-resonance) while minimizing symbolic complexity (binary Sturmian minimality), thus generating a clock texture that is neither periodic nor maximally random.
In the present language, this is a resource-control statement: golden coding minimizes the space overhead needed to avoid resonant collapse of the scan dynamics.

\subsection{Complexity classes: what can and cannot be claimed}
It is tempting to map physical motifs directly onto complexity-class conjectures (e.g. to treat ``persistent matter'' as evidence of class separations such as $\mathrm{P}\ne\mathrm{NP}$).
Such statements are not acceptable as theorem-level claims, because the conjectures are open.

What \emph{is} acceptable is a carefully separated analogy:
interaction can boost verification power in the standard interactive-proof setting \cite{GoldwasserMicaliRackoff1989}, culminating in the theorem $\mathrm{IP}=\mathrm{PSPACE}$ \cite{Shamir1992}.
In an observer-as-interactive-machine picture (Section~\ref{subsec:observer_oracle}), one may view the environment as a prover-like stream and the observer as a verifier-like process constrained by $(T,S)$.
This does not prove any new separation, but it clarifies how ``verification'' and ``construction'' may have different resource profiles in a scan--readout universe.

\subsection{A cybernetic reading}
In cybernetic terms, the observer is a controller operating under resource constraints.
The scan supplies an access channel; the projection supplies a compression channel; undecidability supplies a hard boundary on prediction; and the control objective is to keep the system in a regime where prediction is neither trivial nor impossible.
This yields a concrete, publication-friendly version of ``teleology'': not metaphysical purpose, but resource-stabilized open-ended operation.
